\documentclass[10pt,a4paper]{amsart}
\usepackage[margin=19mm]{geometry}
\usepackage{amsmath,amssymb,mathtools}
\usepackage[scr=rsfs,cal=euler]{mathalfa}
\usepackage{xfrac}
\usepackage[unicode,hidelinks,bookmarks=false]{hyperref}
\newcommand{\ie}{i.e.\ }
\newcommand{\cf}{cf.\ }
\newcommand{\resp}{respectively\ }
\newcommand{\Subsection}{\subsection}
\providecommand{\nobreakdash}{\nobreak\hbox{-}}
\setlength{\emergencystretch}{2em}
\multlinegap0pt
\usepackage{amssymb}
\usepackage{xcolor}
\usepackage{float}
\newtheorem{theorem}{Theorem}
\title{Repetition Avoidance in Curling-Number Transforms}
\author{Geoffrey Caveney, Haoxuan (Jason) Dong, Jeffrey Shallit}
\date{Source: arXiv:2608.15670v1 (CC BY 4.0)}
\begin{document}
\maketitle

\noindent\emph{Source and licence.} Repetition Avoidance in Curling-Number Transforms. Version arXiv:2608.15670v1 (CC BY 4.0), distributed by arXiv under the Creative Commons Attribution 4.0 International licence, http://creativecommons.org/licenses/by/4.0/, as recorded in the arXiv licence metadata for this exact version and shown on the abstract page https://arxiv.org/abs/2608.15670v1. Full licence notice: \texttt{licenses/arxiv-2608.15670-CC-BY-4.0.txt}.\par\smallskip

\noindent\textit{Editorial scope.} Counted unit: the article's theorem thm:ternary-overlap-max (source0.tex lines 286--289). The article's complete original proof of it is delivered with it (source0.tex lines 291--305, one \texttt{proof} environment of 15 lines). No mathematics is written, repaired or re-derived: the statement and the proof are the article's own lines, in the article's own notation, and no retained line is substituted or re-flowed.  No further source-local context is needed: every cross-reference of every retained line resolves inside the two delivered spans.  Nothing was omitted from any retained span, so the package carries no omission record.  No retained line contains a citation, so the wrapper carries no bibliography and no bibliography entry.  No retained line contains a figure environment, an image-inclusion command or drawing code, so no image is delivered and the package declares no asset dependency.  Editorial formatting only, in addition: an AMS \texttt{amsart} preamble that restates the article's own declarations and macros used by the retained lines, namely the environments \texttt{theorem} on separate counters, together with the standard packages those lines need.  The retained environments print in this document as Theorem 1 in order of appearance, whereas the article prints the same items with the numbering of its own full document.  The complete native source of the article as distributed in the arXiv e-print for this exact version is bundled unchanged as \texttt{sources/207771/source0.tex}.
\begin{theorem} 
The longest overlap-free ternary word whose curling-number transform is overlap-free has length $84$.
\label{thm:ternary-overlap-max}
\end{theorem}

\begin{proof} 
The exhaustive search finds $6048$ valid words of length $84$ and none of length $85$, so the maximum length is $84$.  One length-$84$ example, grouped for readability, is
\begin{center}
\small\ttfamily
01101001020020\;21211002002100\;21101001020020\\
21101102002021\;21100200210021\;12201122021100
\end{center}
Its curling-number transform is
\begin{center}
\small\ttfamily
11211221211212\;21122121122112\;12211221211212\\
21211221121221\;12212112211212\;21211212211212
\end{center}
and is also overlap-free.
\end{proof}

\end{document}
